% Options for packages loaded elsewhere
\PassOptionsToPackage{unicode}{hyperref}
\PassOptionsToPackage{hyphens}{url}
\documentclass[
  ignorenonframetext,
]{beamer}
\newif\ifbibliography
\usepackage{pgfpages}
\setbeamertemplate{caption}[numbered]
\setbeamertemplate{caption label separator}{: }
\setbeamercolor{caption name}{fg=normal text.fg}
\beamertemplatenavigationsymbolsempty
% remove section numbering
\setbeamertemplate{part page}{
  \centering
  \begin{beamercolorbox}[sep=16pt,center]{part title}
    \usebeamerfont{part title}\insertpart\par
  \end{beamercolorbox}
}
\setbeamertemplate{section page}{
  \centering
  \begin{beamercolorbox}[sep=12pt,center]{section title}
    \usebeamerfont{section title}\insertsection\par
  \end{beamercolorbox}
}
\setbeamertemplate{subsection page}{
  \centering
  \begin{beamercolorbox}[sep=8pt,center]{subsection title}
    \usebeamerfont{subsection title}\insertsubsection\par
  \end{beamercolorbox}
}
% Prevent slide breaks in the middle of a paragraph
\widowpenalties 1 10000
\raggedbottom
\AtBeginPart{
  \frame{\partpage}
}
\AtBeginSection{
  \ifbibliography
  \else
    \frame{\sectionpage}
  \fi
}
\AtBeginSubsection{
  \frame{\subsectionpage}
}
\usepackage{iftex}
\ifPDFTeX
  \usepackage[T1]{fontenc}
  \usepackage[utf8]{inputenc}
  \usepackage{textcomp} % provide euro and other symbols
\else % if luatex or xetex
  \usepackage{unicode-math} % this also loads fontspec
  \defaultfontfeatures{Scale=MatchLowercase}
  \defaultfontfeatures[\rmfamily]{Ligatures=TeX,Scale=1}
\fi
\usepackage{lmodern}
\ifPDFTeX\else
  % xetex/luatex font selection
\fi
% Use upquote if available, for straight quotes in verbatim environments
\IfFileExists{upquote.sty}{\usepackage{upquote}}{}
\IfFileExists{microtype.sty}{% use microtype if available
  \usepackage[]{microtype}
  \UseMicrotypeSet[protrusion]{basicmath} % disable protrusion for tt fonts
}{}
\makeatletter
\@ifundefined{KOMAClassName}{% if non-KOMA class
  \IfFileExists{parskip.sty}{%
    \usepackage{parskip}
  }{% else
    \setlength{\parindent}{0pt}
    \setlength{\parskip}{6pt plus 2pt minus 1pt}}
}{% if KOMA class
  \KOMAoptions{parskip=half}}
\makeatother
\setlength{\emergencystretch}{3em} % prevent overfull lines
\providecommand{\tightlist}{%
  \setlength{\itemsep}{0pt}\setlength{\parskip}{0pt}}
% Slide layout defaults for warning-clean scientific decks.
\usepackage{etoolbox}
\IfFileExists{xurl.sty}{\usepackage{xurl}}{}
\IfFileExists{seqsplit.sty}{\usepackage{seqsplit}}{\newcommand{\seqsplit}[1]{#1}}
\protected\def\breakseq#1{\seqsplit{#1}}
\protected\def\breaktt#1{\begingroup\ttfamily\seqsplit{#1}\endgroup}
\setlength{\emergencystretch}{6em}
\tolerance=5000
\hbadness=10000
\hfuzz=1pt
\setlength{\tabcolsep}{2pt}
\AtBeginEnvironment{longtable}{\tiny\renewcommand{\arraystretch}{0.86}\setlength{\tabcolsep}{1pt}}
\AtBeginEnvironment{tabular}{\tiny\renewcommand{\arraystretch}{0.86}\setlength{\tabcolsep}{1pt}}

% Natbib and cross-reference fallbacks — slides don't load natbib
% or cleveref, but combined-PDF manuscript prose may emit these
% commands. The fallback renders citations as a bracketed key list
% and cross-references as detokenized labels so slides don't fail on
% undefined control sequences or raw underscores. \providecommand is
% a no-op if packages load later.
\providecommand{\citep}[1]{[#1]}
\providecommand{\citet}[1]{#1}
\providecommand{\citealp}[1]{#1}
\providecommand{\citeauthor}[1]{#1}
\providecommand{\citeyear}[1]{#1}
\providecommand{\cref}[1]{\texttt{\detokenize{#1}}}
\providecommand{\Cref}[1]{\texttt{\detokenize{#1}}}
\usepackage{bookmark}
\IfFileExists{xurl.sty}{\usepackage{xurl}}{} % add URL line breaks if available
\urlstyle{same}
% fallback for those not using the hyperref driver hyperxmp:
\makeatletter
\@ifundefined{xmpquote}{\newcommand{\xmpquote}[1]{#1}}{}
\makeatother
\hypersetup{
  hidelinks,
  pdfcreator={LaTeX via pandoc}}

\author{\texorpdfstring{}{}}
\date{}

\begin{document}

\section{Discussion}\label{discussion}

\begin{frame}[fragile,allowframebreaks]{The Zero-Mock Tradeoff}
\protect\phantomsection\label{the-zero-mock-tradeoff}
The \href{03e_quality.md\#zero-mock-testing-policy}{Zero-Mock testing
policy} is \texttt{template/}'s most distinctive design decision. By
prohibiting all mock objects, we gain confidence that tests exercise
real code paths---a pytest run against the template genuinely invokes
\texttt{pandoc}, writes to disk, and parses real YAML. The cost is test
duration: the full infrastructure test suite (\textasciitilde7,968
tests) takes 2--4 minutes, compared to sub-second execution typical of
heavily-mocked suites.

We argue this tradeoff is strongly favorable for research software.
Unlike web applications where millisecond latency and thousands of daily
deploys demand fast feedback loops, research pipelines run infrequently
(once per manuscript revision) and correctness vastly outweighs speed. A
mocked test that passes while the real renderer fails is worse than a
slow test that catches the failure. The analogy to statistical
methodology is precise: just as Simmons et al.'s \emph{researcher
degrees of freedom} \breakseq{[@simmons2011falsepositive]} inflate
false-positive rates through undisclosed analytical flexibility, mock
objects create \emph{testing degrees of freedom} that make integration
failures invisible. The Zero-Mock policy closes this loophole by the
same mechanism that pre-registration \breakseq{[@nosek2018preregistration]}
closes the p-hacking loophole: removing flexibility before the fact. As
Peng \breakseq{[@peng2011reproducible]} argues, computational reproducibility
requires independent verification---and mock-only tests verify
assumptions rather than results. Garijo et al.'s FAIRsoft evaluator
\breakseq{[@garijo2024fairsoft]} identifies \emph{executability} as a primary
quality indicator; the Zero-Mock policy operationalizes executability at
the unit level.

\begin{block}{When Mocks Are Not the Problem}
\protect\phantomsection\label{when-mocks-are-not-the-problem}
It is important to distinguish the Zero-Mock policy from a naive
rejection of all test isolation techniques. Fowler's classification
\breakseq{[@martin2008clean]} recognizes that stubs and fakes serve legitimate
purposes---a test database populated with known data is not a mock, it
is a fixture. The policy specifically prohibits \emph{mock objects} as
defined by Meszaros: assertions on indirect outputs (method calls,
argument patterns) rather than direct outputs (return values, side
effects). The distinction matters because mock-based assertions encode
implementation assumptions (``method X must be called with argument Y'')
that become invisible coupling between tests and production code,
creating the illusion of coverage without testing real behavior.
\end{block}

\begin{block}{Practical Implementation}
\protect\phantomsection\label{practical-implementation}
The template enforces zero-mock compliance at three levels:

\begin{enumerate}
\tightlist
\item
  \textbf{Code review}: \texttt{AGENTS.md} at every directory level
  explicitly states the prohibition, ensuring both human and AI
  contributors are aware before writing tests.
\item
  \textbf{Static analysis}:
  \texttt{grep\ -rn\ "MagicMock\textbackslash{}\textbar{}unittest.mock\textbackslash{}\textbar{}@patch"\ tests/}
  can be run as a pre-commit hook to catch violations.
\item
  \textbf{Cultural norm}: \breaktt{template\_code\_project} documents
  filesystem + YAML + plotting paths while
  \breaktt{template\_autoresearch\_project} exercises readiness planning;
  \breaktt{template\_search\_project} reinforces HTTP-realistic
  literature queries---both serve as onboarding references alongside
  infra suites.
\end{enumerate}

However, the policy requires careful management of external
dependencies. Tests requiring Ollama (the local LLM backend) use
\breaktt{@pytest.mark.requires\_ollama} and are skipped in environments
where the service is unavailable. Tests requiring network access use
\breaktt{@pytest.mark.network}. This marker system preserves the
Zero-Mock principle while acknowledging that not all environments
provide all services, especially computationally intensive ones. The key
distinction is between \emph{replacing} an external dependency (which
mock objects do, hiding failures) and \emph{skipping} a test when a
dependency is absent (which markers do, preserving transparency).
\end{block}
\end{frame}

\end{document}
